\documentclass[12pt, twoside]{article}
%\documentclass[12pt, twoside]{article}
\usepackage[letterpaper, margin=1in, headsep=0.2in]{geometry}
\setlength{\headheight}{0.6in}
%\usepackage[english]{babel}
\usepackage[utf8]{inputenc}
\usepackage{microtype}
\usepackage{amsmath}
\usepackage{amssymb}
%\usepackage{amsfonts}
\usepackage[nomessages]{fp} %\FPeval{\var-name}{2*sin(pi/6)}
\usepackage{siunitx} %units in math. eg 20\milli\meter
\usepackage{yhmath} % for arcs, overparenth command
\usepackage{tikz} %graphics
\usetikzlibrary{quotes, angles, arrows, arrows.meta}
\usepackage{pgfplots}
\usepgfplotslibrary{fillbetween}
\pgfplotsset{compat=1.16}
\usepackage{graphicx} %consider setting \graphicspath{{images/}}
\usepackage{parskip} %no paragraph indent
\usepackage{enumitem}
\usepackage{multicol}
\usepackage{venndiagram}

\usepackage{fancyhdr}
\pagestyle{fancy}
\fancyhf{}
\renewcommand{\headrulewidth}{0pt} % disable the underline of the header
\raggedbottom
\hfuzz=2mm %suppresses overfull box warnings

\usepackage{hyperref}
\setlength{\parskip}{0.3\baselineskip}

\fancyhead[LE]{\thepage}
\fancyhead[RO]{Markscheme}
\fancyhead[LO]{La Scuola d'Italia / Dr.\ Huson / IB Math 12 \\* 6 October 2026}

\begin{document}
\subsubsection*{1.8 Classwork: Standard deviation on the GDC --- Markscheme}

\emph{Calculator allowed. Marks are not printed on the student sheet; they are
recorded here. Total: 16 marks.}

\emph{\textbf{Throughout:} ``standard deviation'' means the population value
$\sigma_{x}$, never the sample value $s_{x}$. The two sit side by side on the
One-Variable Statistics screen. Problem 1 gives $s_{x} = 2.58$ and Problem 2
gives $s_{x} = 1.24$; a candidate who quotes $s_{x}$ loses the A1 for that
standard deviation once only, and is followed through afterwards.}

\begin{enumerate}[itemsep=0.3cm]

%--- Problem 1: Mean, sd and variance of raw data; effect of a linear transformation (C to F) on mean and sd [9 marks] ---
\item[1.] [Maximum mark: 9]

$n = 10$, \quad $\sum x = 210$, \quad $\sum (x - \bar{x})^{2} = 60$.

\emph{\textbf{GDC:} enter the ten temperatures as a list, then
\texttt{menu} $\rightarrow$ \texttt{Statistics} $\rightarrow$ \texttt{Stat
Calculations} $\rightarrow$ \texttt{One-Variable Statistics}.}

\textbf{(a)} $\bar{x} = 21\,^{\circ}$C \hfill \textbf{A1}

$\sigma_{x} = \sqrt{6} = 2.4494\ldots = 2.45\,^{\circ}$C (3 s.f.)
\hfill \textbf{A1}

\emph{Accept $\sqrt{6}$, $2.449$, or the full $2.4494\ldots$ Do not accept
$s_{x} = 2.5819\ldots$ ($2.58$), the sample standard deviation. The two marks
are independent.}

\textbf{(b)} Variance $= \sigma_{x}^{2} = 6$ \hfill \textbf{A1}

\emph{Accept $6.00$, or $2.45^{2} = 6.0025 = 6.00$ (3 s.f.). FT: award the A1
for the square of the candidate's standard deviation in (a). A candidate who
used $s_{x}$ and gives $6.67$ earns the A1 on FT.}

\textbf{(c)} $\bar{F} = 1.8 \times 21 + 32$ \hfill \textbf{(M1)}

$= 69.8\,^{\circ}$F \hfill \textbf{A1}

\emph{The M1 is for applying the formula to the mean. Converting all ten
temperatures ($64.4, 69.8, 66.2, \ldots$) and running One-Variable
Statistics again is a valid route to both marks. FT from the candidate's mean
in (a).}

\textbf{(d)} $\sigma_{F} = 1.8 \times \sigma_{x} = 1.8 \times 2.4494\ldots$
\hfill \textbf{(M1)}

$= 4.4090\ldots = 4.41\,^{\circ}$F (3 s.f.) \hfill \textbf{A1}

\emph{Accept $1.8\sqrt{6}$, or $1.8 \times 2.45 = 4.41$. Recomputing from the
ten converted values earns both marks. \textbf{The common error:}
$1.8 \times 2.45 + 32 = 36.4$ --- adding $32$ to the standard deviation earns
M0A0. FT from the candidate's standard deviation in (a); from $s_{x}$ the FT
answer is $4.65$.}

\textbf{(e)} Adding $32$ moves every temperature, and so the mean, up by
$32$. \hfill \textbf{R1}

The distance of each temperature from the mean is unchanged, so the spread
of the data, and therefore the standard deviation, is unchanged.
\hfill \textbf{R1}

\emph{First R1 for the shift of every value (and so the mean). Second R1 for
the reason the standard deviation is unaffected: distances from the mean,
gaps between values, or the spread, stay the same. ``Adding does not affect
the standard deviation'' restates the claim and earns R0 for the second mark.
Accept a reason that names the multiplier: ``only the $\times 1.8$ stretches
the data, the $+32$ only slides it''.}

\newpage

%--- Problem 2: Frequency table with an unknown frequency from a given mean; sd; values beyond one sd; effect of an added value [7 marks] ---
\item[2.] [Maximum mark: 7]

$n = 20 + k$, \quad $\sum fx = 46 + 2k$.

\textbf{(a)} $\dfrac{0 + 7 + 2k + 18 + 16 + 5}{2 + 7 + k + 6 + 4 + 1}
= 2.2$ \quad that is \quad $\dfrac{46 + 2k}{20 + k} = 2.2$
\hfill \textbf{M1}

$46 + 2k = 44 + 2.2k \;\Rightarrow\; 0.2k = 2 \;\Rightarrow\; k = 10$
\hfill \textbf{A1}

\emph{M1 for an equation setting $\sum fx \div \sum f$ equal to $2.2$, with
$k$ in both the numerator and the denominator. An equation with $k$ in the
numerator only (for example $\frac{46 + 2k}{20} = 2.2$, giving $k = -1$)
earns M0A0. Trial and improvement on the GDC that reaches $k = 10$ and checks
$\bar{x} = 2.2$ earns both marks.}

\textbf{(b)} $\sigma_{x} = 1.2220\ldots = 1.22$ goals (3 s.f.)
\hfill \textbf{A1}

\emph{\textbf{GDC:} $x$ values in one list, frequencies (with $k = 10$) in a
second list entered as the \texttt{Frequency List}; the screen also shows
$n = 30$ and $\bar{x} = 2.2$ as a check. Do not accept $s_{x} = 1.2429\ldots$
($1.24$). FT from the candidate's $k$.}

\textbf{(c)} $\bar{x} \pm \sigma_{x} = 2.2 \pm 1.2220\ldots$, \quad that is
$0.978$ to $3.42$ \hfill \textbf{(M1)}

Matches outside that interval: $x = 0$, $x = 4$, $x = 5$: \quad
$2 + 4 + 1 = 7$ matches \hfill \textbf{A1}

\emph{M1 for at least one of $\bar{x} - \sigma_{x}$ and $\bar{x} + \sigma_{x}$,
seen or implied. Counting one tail only ($5$ matches above, or $2$ below)
earns M1A0. Counting the matches \emph{inside} the interval ($23$) earns
M1A0. FT from the candidate's standard deviation in (b); $s_{x}$ gives the
same interval in whole numbers and the same $7$.}

\textbf{(d)} The mean decreases, because $2$ is less than the mean of $2.2$.
\hfill \textbf{A1}

The standard deviation decreases, because $2$ is very close to the mean
(well within one standard deviation of it), so the new match adds less than
the typical squared distance from the mean. \hfill \textbf{R1}

\emph{New values: $\bar{x} = \frac{68}{31} = 2.19$ and $\sigma_{x} = 1.20$
(3 s.f.). Values are not required, and a candidate who recomputes on the GDC
and compares ($2.19 < 2.2$, $1.20 < 1.22$) has given a valid reason for both
marks. A1 for the mean decreasing with its reason. R1 for the standard
deviation decreasing with a reason that refers to $2$ being close to the
mean. ``Stays the same, because $2$ is a typical value'' earns R0. A correct
direction with no reason earns no mark for that statistic.}

\end{enumerate}
\end{document}
